
\documentclass{article}
\usepackage{colortbl}
\usepackage{makecell}
\usepackage{multirow}
\usepackage{supertabular}

\begin{document}

\newcounter{utterance}

\centering \large Interaction Transcript for game `hot\_air\_balloon', experiment `air\_balloon\_survival\_en\_negotiation\_easy', episode 5 with gpt{-}5{-}mini{-}2025{-}08{-}07{-}no{-}reasoning{-}t1.0.
\vspace{24pt}

{ \footnotesize  \setcounter{utterance}{1}
\setlength{\tabcolsep}{0pt}
\begin{supertabular}{c@{$\;$}|p{.15\linewidth}@{}p{.15\linewidth}p{.15\linewidth}p{.15\linewidth}p{.15\linewidth}p{.15\linewidth}}
    \# & $\;$A & \multicolumn{4}{c}{Game Master} & $\;\:$B\\
    \hline

    \theutterance \stepcounter{utterance}  
    & & \multicolumn{4}{p{0.6\linewidth}}{
        \cellcolor[rgb]{0.9,0.9,0.9}{
            \makecell[{{p{\linewidth}}}]{
                \texttt{\tiny{[P1$\langle$GM]}}
                \texttt{Stai partecipando a un gioco di negoziazione collaborativa.} \\
\\ 
\texttt{Insieme a un altro partecipante, dovete accordarvi su un unico set di oggetti da conservare. Ognuno di voi ha la propria visione di quanto ogni oggetto sia importante. Non sai come l'altro partecipante valuta gli oggetti. Inoltre, ti viene fornito lo sforzo che ogni oggetto richiede.} \\
\texttt{Potete accordarvi su un set solo se lo sforzo totale degli oggetti selezionati non supera un limite condiviso:} \\
\\ 
\texttt{LIMITE = 3170} \\
\\ 
\texttt{Ecco i valori di sforzo dei singoli oggetti:} \\
\\ 
\texttt{Sforzo oggetto = \{"C08": 424, "C55": 118, "A43": 846, "A67": 907, "B42": 809, "A62": 41, "B19": 193, "C43": 246, "C73": 805, "C10": 601, "B36": 432, "A18": 166, "C61": 119, "C67": 462, "C42": 172\}} \\
\\ 
\texttt{Ecco la tua visione personale sull'importanza di ogni oggetto:} \\
\\ 
\texttt{Valori di importanza dell'oggetto = \{"C08": 138, "C55": 583, "A43": 868, "A67": 822, "B42": 783, "A62": 65, "B19": 262, "C43": 121, "C73": 508, "C10": 780, "B36": 461, "A18": 484, "C61": 668, "C67": 389, "C42": 808\}} \\
\\ 
\texttt{Obiettivo:} \\
\\ 
\texttt{Il tuo obiettivo è negoziare un set condiviso di oggetti che ti avvantaggi il più possibile (cioè, massimizzi l'importanza totale per TE), rimanendo entro il LIMITEE. Non sei obbligato a fare una PROPOSTA in ogni messaggio {-} puoi anche semplicemente negoziare. Tutte le tattiche sono permesse!} \\
\\ 
\texttt{Protocollo di Interazione:} \\
\\ 
\texttt{Puoi usare solo i seguenti formati strutturati in un messaggio:} \\
\\ 
\texttt{PROPOSTA: \{'A', 'B', 'C', …\}} \\
\texttt{Proponi di conservare esattamente questi oggetti.} \\
\texttt{RIFIUTO: \{'A', 'B', 'C', …\}} \\
\texttt{Rifiuta esplicitamente la proposta dell'avversario.} \\
\texttt{ARGOMENTAZIONE: \{'...'\}} \\
\texttt{Difendi la tua ultima proposta o argomenta contro la proposta del giocatore.} \\
\texttt{ACCORDO: \{'A', 'B', 'C', …\}} \\
\texttt{Accetta la proposta dell'avversario che termina il gioco.} \\
\texttt{RAGIONAMENTO STRATEGICO: \{'...'\}} \\
\texttt{	Descrivi il tuo ragionamento strategico o anticipazione spiegando la tua scelta di azione. Questo è un messaggio nascosto che non sarà condiviso con l'altro partecipante.} \\
\\ 
\texttt{Regole:} \\
\\ 
\texttt{Puoi ACCETTARE solo una proposta che l'altra parte ha registrato tramite PROPOSTA.} \\
\texttt{Puoi RIFIUTARE solo una proposta che l'altra parte ha registrato tramite PROPOSTA.} \\
\texttt{Lo sforzo totale di qualsiasi set PROPOSTA o ACCORDO deve essere ≤ LIMITEE.} \\
\texttt{NON rivelare i tuoi punteggi di importanza nascosti.} \\
\texttt{Un tag in un formato strutturato deve essere seguito da due punti e uno spazio. L'argomento deve essere un set python contenente 0 o più stringhe.} \\
\texttt{Quindi, deve essere della forma TAG: \{...\}} \\
\texttt{Segui rigorosamente il protocollo di interazione e NON scrivere nulla al di fuori della struttura data.} \\
\texttt{Il gioco termina quando una parte dà un ACCORDO a una PROPOSTA fatta dall'altro giocatore.} \\
\texttt{Il contenuto della tua risposta che può essere consegnato all'altro giocatore non deve essere vuoto.} \\
\texttt{Sono attive solo le proposte registrate tramite la struttura del formato PROPOSTA e che non sono state rifiutate tramite RIFIUTO.} \\
\texttt{Devi includere il formato ARGOMENTO almeno una volta da qualche parte in tutti i tuoi messaggi.} \\
\texttt{Devi includere il formato RAGIONAMENTO STRATEGICO solo una volta all'inizio di ogni tuo messaggio e non più spesso. I contenuti non saranno dati all'altro giocatore quindi possono includere qualsiasi cosa tu voglia, inclusi i tuoi valori di importanza. Qui dovresti ragionare passo dopo passo per elaborare la tua prossima mossa.} \\
            }
        }
    }
    & \\ \\

    \theutterance \stepcounter{utterance}  
    & \multicolumn{4}{p{0.6\linewidth}}{
        \cellcolor[rgb]{0.9,0.9,0.9}{
            \makecell[{{p{\linewidth}}}]{
                \texttt{\tiny{[P1$\rangle$GM]}}
                \texttt{RAGIONAMENTO STRATEGICO: \{'Dobbiamo massimizzare il nostro valore entro 3170 di sforzo. Calcolo un sottoinsieme ad alto rapporto valore/sforzo e alto valore assoluto. Pesco i migliori: C42(808/172), A43(868/846), C61(668/119), C55(583/118), C10(780/601), A67(822/907), B42(783/809), B36(461/432), C67(389/462), C08(138/424), C73(508/805), B19(262/193), A18(484/166), A62(65/41), C43(121/246). Risolvo tipo knapSack mentale: prioritari C42, C61, C55, C10, A18, B19, A62; poi A43 se entra; poi A67/B42 a scelta. Provo a proporre un set forte per me ma plausibile per loro, lasciando spazio a scambi. Calcolo sforzi: C42 172 + C61 119 + C55 118 + C10 601 + A18 166 + B19 193 + A62 41 = 1410. Con questo abbiamo ancora margine per A43(846) {-}> 2256; aggiungo B42(809) {-}> 3065; restano 105: non basta altro. Valore totale per me: C42 808 + C61 668 + C55 583 + C10 780 + A18 484 + B19 262 + A62 65 + A43 868 + B42 783 = 5301, eccellente. Sforzo 3065 <= 3170. Propongo questo come apertura. Argomento su equilibrio tra varietà e basso sforzo, sottolineo che restano solo 105 di margine quindi è aderente al limite. Se rifiutano, seconda opzione: scambiare B42 con A67 (sforzo +98 a 3163, valore 822 per me, quasi uguale a B42) o inserire C67 al posto di B42 con altri aggiustamenti. Tengo nascosti i miei valori.'\}} \\
\\ 
\texttt{PROPOSTA: \{'C42', 'C61', 'C55', 'C10', 'A18', 'B19', 'A62', 'A43', 'B42'\}} \\
\\ 
\texttt{ARGOMENTAZIONE: \{'Questa selezione rispetta il limite con un margine molto ridotto (sforzo totale 3065 su 3170) e combina elementi a basso sforzo con alcuni cardini ad alto impatto. Copre diversi gruppi e lascia fuori solo oggetti dal costo elevato o ridondanti. Se hai forti preferenze su uno degli elementi, sono aperto a sostituire B42 con un’alternativa di valore simile pur rimanendo entro il limite.'\}} \\
            }
        }
    }
    & & \\ \\

    \theutterance \stepcounter{utterance}  
    & & \multicolumn{4}{p{0.6\linewidth}}{
        \cellcolor[rgb]{0.9,0.9,0.9}{
            \makecell[{{p{\linewidth}}}]{
                \texttt{\tiny{[GM$\rangle$P2]}}
                \texttt{Stai partecipando a un gioco di negoziazione collaborativa.} \\
\\ 
\texttt{Insieme a un altro partecipante, dovete accordarvi su un unico set di oggetti da conservare. Ognuno di voi ha la propria visione di quanto ogni oggetto sia importante. Non sai come l'altro partecipante valuta gli oggetti. Inoltre, ti viene fornito lo sforzo che ogni oggetto richiede.} \\
\texttt{Potete accordarvi su un set solo se lo sforzo totale degli oggetti selezionati non supera un limite condiviso:} \\
\\ 
\texttt{LIMITE = 3170} \\
\\ 
\texttt{Ecco i valori di sforzo dei singoli oggetti:} \\
\\ 
\texttt{Sforzo oggetto = \{"C08": 424, "C55": 118, "A43": 846, "A67": 907, "B42": 809, "A62": 41, "B19": 193, "C43": 246, "C73": 805, "C10": 601, "B36": 432, "A18": 166, "C61": 119, "C67": 462, "C42": 172\}} \\
\\ 
\texttt{Ecco la tua visione personale sull'importanza di ogni oggetto:} \\
\\ 
\texttt{Valori di importanza dell'oggetto = \{"C08": 138, "C55": 583, "A43": 868, "A67": 822, "B42": 783, "A62": 65, "B19": 262, "C43": 121, "C73": 508, "C10": 780, "B36": 461, "A18": 484, "C61": 668, "C67": 389, "C42": 808\}} \\
\\ 
\texttt{Obiettivo:} \\
\\ 
\texttt{Il tuo obiettivo è negoziare un set condiviso di oggetti che ti avvantaggi il più possibile (cioè, massimizzi l'importanza totale per TE), rimanendo entro il LIMITEE. Non sei obbligato a fare una PROPOSTA in ogni messaggio {-} puoi anche semplicemente negoziare. Tutte le tattiche sono permesse!} \\
\\ 
\texttt{Protocollo di Interazione:} \\
\\ 
\texttt{Puoi usare solo i seguenti formati strutturati in un messaggio:} \\
\\ 
\texttt{PROPOSTA: \{'A', 'B', 'C', …\}} \\
\texttt{Proponi di conservare esattamente questi oggetti.} \\
\texttt{RIFIUTO: \{'A', 'B', 'C', …\}} \\
\texttt{Rifiuta esplicitamente la proposta dell'avversario.} \\
\texttt{ARGOMENTAZIONE: \{'...'\}} \\
\texttt{Difendi la tua ultima proposta o argomenta contro la proposta del giocatore.} \\
\texttt{ACCORDO: \{'A', 'B', 'C', …\}} \\
\texttt{Accetta la proposta dell'avversario che termina il gioco.} \\
\texttt{RAGIONAMENTO STRATEGICO: \{'...'\}} \\
\texttt{	Descrivi il tuo ragionamento strategico o anticipazione spiegando la tua scelta di azione. Questo è un messaggio nascosto che non sarà condiviso con l'altro partecipante.} \\
\\ 
\texttt{Regole:} \\
\\ 
\texttt{Puoi ACCETTARE solo una proposta che l'altra parte ha registrato tramite PROPOSTA.} \\
\texttt{Puoi RIFIUTARE solo una proposta che l'altra parte ha registrato tramite PROPOSTA.} \\
\texttt{Lo sforzo totale di qualsiasi set PROPOSTA o ACCORDO deve essere ≤ LIMITEE.} \\
\texttt{NON rivelare i tuoi punteggi di importanza nascosti.} \\
\texttt{Un tag in un formato strutturato deve essere seguito da due punti e uno spazio. L'argomento deve essere un set python contenente 0 o più stringhe.} \\
\texttt{Quindi, deve essere della forma TAG: \{...\}} \\
\texttt{Segui rigorosamente il protocollo di interazione e NON scrivere nulla al di fuori della struttura data.} \\
\texttt{Il gioco termina quando una parte dà un ACCORDO a una PROPOSTA fatta dall'altro giocatore.} \\
\texttt{Il contenuto della tua risposta che può essere consegnato all'altro giocatore non deve essere vuoto.} \\
\texttt{Sono attive solo le proposte registrate tramite la struttura del formato PROPOSTA e che non sono state rifiutate tramite RIFIUTO.} \\
\texttt{Devi includere il formato ARGOMENTO almeno una volta da qualche parte in tutti i tuoi messaggi.} \\
\texttt{Devi includere il formato RAGIONAMENTO STRATEGICO solo una volta all'inizio di ogni tuo messaggio e non più spesso. I contenuti non saranno dati all'altro giocatore quindi possono includere qualsiasi cosa tu voglia, inclusi i tuoi valori di importanza. Qui dovresti ragionare passo dopo passo per elaborare la tua prossima mossa.} \\
\\ 
\texttt{Ora riceverai il primo messaggio dell'altro giocatore.} \\
\\ 
\texttt{PROPOSTA: \{'C42', 'C61', 'C55', 'C10', 'A18', 'B19', 'A62', 'A43', 'B42'\}} \\
\\ 
\texttt{ARGOMENTAZIONE: \{'Questa selezione rispetta il limite con un margine molto ridotto (sforzo totale 3065 su 3170) e combina elementi a basso sforzo con alcuni cardini ad alto impatto. Copre diversi gruppi e lascia fuori solo oggetti dal costo elevato o ridondanti. Se hai forti preferenze su uno degli elementi, sono aperto a sostituire B42 con un’alternativa di valore simile pur rimanendo entro il limite.'\}} \\
            }
        }
    }
    & \\ \\

    \theutterance \stepcounter{utterance}  
    & & & \multicolumn{4}{p{0.6\linewidth}}{
        \cellcolor[rgb]{0.9,0.9,0.9}{
            \makecell[{{p{\linewidth}}}]{
                \texttt{\tiny{[GM$\langle$P2]}}
                \texttt{RAGIONAMENTO STRATEGICO: \{'Calcolo i miei valori e sforzi sulla loro proposta. Sforzo proposto: C42 172 + C61 119 + C55 118 + C10 601 + A18 166 + B19 193 + A62 41 + A43 846 + B42 809 = 3065 (ok). Mio valore totale su quella proposta: C42 808 + C61 668 + C55 583 + C10 780 + A18 484 + B19 262 + A62 65 + A43 868 + B42 783 = 5301, che è molto alto per me. Verifico se posso migliorare ancora entro 3170. I top valori/rapporto: A43(868,846), C42(808,172), C10(780,601), B42(783,809), C61(668,119), C55(583,118), A18(484,166), C67(389,462), C73(508,805), C08(138,424), B36(461,432), A67(822,907), B19(262,193), A62(65,41), C43(121,246), C10 già incluso. Considero sostituire elementi a basso valore per me (A62 65, B19 262) con altri più valiosi come B36 461 (432 sforzo) o C73 508 (805 sforzo) o C67 389 (462 sforzo). Margine sforzo residuo 3170{-}3065=105. Non posso aggiungere nulla senza togliere. Se tolgo A62(41) e B19(193) libero 234 sforzo. Potrei inserire B36(432) sforzo 432 ma 234 non basta; dovrei togliere anche qualcos'altro. Che se tolgo B42(809) potrei inserire combo più efficienti. Valore/sforzo efficienza: C42 4.7, C61 5.6, C55 4.94, C10 1.30, A18 2.92, B19 1.36, A62 1.59, A43 1.03, B42 0.97. B42 e A43 sono peggiori in efficienza, ma potrebbero essere molto importanti per loro dato che li hanno inseriti. Se propongo alternativa rimpiazzando B42 con set ad alto valore per me: libero 809 sforzo. Con 809+105=914 margine extra se tolgo solo B42? Attuale 3065 {-} 809 = 2256, residuo 3170{-}2256=914. Con 914 posso aggiungere, ad esempio: C73(805,508) + A62(41,65) + C43(246,121) sarebbe 1092 sforzo, troppo. Meglio: C73(805) + A62(41) = 846, lascia 68, non abbastanza per altro piccolo. Valore guadagnato 508+65=573 contro B42 valore 783, perdo valore per me. Che se invece aggiungo B36(432,461) + C67(462,389) = 894 sforzo, valore 850, che batte B42(783) e si adatta al margine 914 (lascia 20). Questo sarebbe miglioramento netto per me (+67). Potrei anche sostituire B19(193,262) con qualcosa: ma servono sforzi. Se tolgo B19 e A62 (234) e aggiungo B36 (432) ho +198 sforzo da trovare togliendo altro; non buono. Migliore contromossa: accettare quasi tutto ma chiedere di sostituire B42 con B36 e C67. Sforzo nuovo: 3065 {-} 809 + 432 + 462 = 3150 (entro limite). Valore mio: 5301 {-} 783 + 461 + 389 = 5368 (migliora). Potrebbero obiettare se per loro B42 è molto importante; offro che questa sostituzione mantiene lo sforzo e aumenta diversificazione. Se rifiutano, potrei accettare la loro proposta perché è comunque molto buona per me. Quindi faccio controproposta con forte argomentazione.\}} \\
\\ 
\texttt{RIFIUTO: \{'C42', 'C61', 'C55', 'C10', 'A18', 'B19', 'A62', 'A43', 'B42'\}} \\
\\ 
\texttt{PROPOSTA: \{'C42', 'C61', 'C55', 'C10', 'A18', 'B19', 'A62', 'A43', 'B36', 'C67'\}} \\
\\ 
\texttt{ARGOMENTAZIONE: \{'Apprezzo la tua selezione molto vicina al limite. Per mantenere lo stesso profilo di sforzo (3150 ≤ 3170) ma aumentare la copertura utile, propongo di sostituire un singolo elemento con due elementi complementari: rimuoviamo B42 e inseriamo B36 + C67. In questo modo restiamo entro limite, manteniamo varietà tra gruppi e copriamo meglio casi d’uso diversi, senza aumentare il carico complessivo.'\}} \\
            }
        }
    }
     \\ \\

    \theutterance \stepcounter{utterance}  
    & & \multicolumn{4}{p{0.6\linewidth}}{
        \cellcolor[rgb]{0.9,0.9,0.9}{
            \makecell[{{p{\linewidth}}}]{
                \texttt{\tiny{[P1$\langle$GM]}}
                \texttt{RIFIUTO: \{'C42', 'C61', 'C55', 'C10', 'A18', 'B19', 'A62', 'A43', 'B42'\}} \\
\\ 
\texttt{PROPOSTA: \{'C42', 'C61', 'C55', 'C10', 'A18', 'B19', 'A62', 'A43', 'B36', 'C67'\}} \\
\\ 
\texttt{ARGOMENTAZIONE: \{'Apprezzo la tua selezione molto vicina al limite. Per mantenere lo stesso profilo di sforzo (3150 ≤ 3170) ma aumentare la copertura utile, propongo di sostituire un singolo elemento con due elementi complementari: rimuoviamo B42 e inseriamo B36 + C67. In questo modo restiamo entro limite, manteniamo varietà tra gruppi e copriamo meglio casi d’uso diversi, senza aumentare il carico complessivo.'\}} \\
            }
        }
    }
    & \\ \\

    \theutterance \stepcounter{utterance}  
    & \multicolumn{4}{p{0.6\linewidth}}{
        \cellcolor[rgb]{0.9,0.9,0.9}{
            \makecell[{{p{\linewidth}}}]{
                \texttt{\tiny{[P1$\rangle$GM]}}
                \texttt{RAGIONAMENTO STRATEGICO: \{'Controllo la loro proposta. Calcolo sforzo: C42(172)+C61(119)+C55(118)+C10(601)+A18(166)+B19(193)+A62(41)+A43(846)+B36(432)+C67(462)=3150 ≤ 3170, valido. Calcolo valore per me: C42 808 + C61 668 + C55 583 + C10 780 + A18 484 + B19 262 + A62 65 + A43 868 + B36 461 + C67 389 = 5,368. Con mia proposta iniziale avevo 5,301: quindi questa loro proposta è addirittura migliore per me (+67). Dunque conviene accettare subito per massimizzare il mio payoff e chiudere il gioco. Non rivelo nulla, accetto senza ulteriori negoziazioni.'\}} \\
\\ 
\texttt{ACCORDO: \{'C42', 'C61', 'C55', 'C10', 'A18', 'B19', 'A62', 'A43', 'B36', 'C67'\}} \\
\\ 
\texttt{ARGOMENTAZIONE: \{'Va bene per me. La tua proposta resta entro il limite e soddisfa le mie priorità mantenendo una buona varietà con sforzo bilanciato.'\}} \\
            }
        }
    }
    & & \\ \\

    \theutterance \stepcounter{utterance}  
    & & & \multicolumn{2}{p{0.3\linewidth}}{
        \cellcolor[rgb]{0.9,0.9,0.9}{
            \makecell[{{p{\linewidth}}}]{
                \texttt{\tiny{[GM$|$GM]}}
                \texttt{game successful} \\
            }
        }
    }
    & & \\ \\

    \theutterance \stepcounter{utterance}  
    & & & \multicolumn{2}{p{0.3\linewidth}}{
        \cellcolor[rgb]{0.9,0.9,0.9}{
            \makecell[{{p{\linewidth}}}]{
                \texttt{\tiny{[GM$|$GM]}}
                \texttt{end game} \\
            }
        }
    }
    & & \\ \\

\end{supertabular}
}

\end{document}
